\documentclass[11pt]{amsart}
\usepackage{mathmemo}
\renewcommand{\memotitle}{OPAC-012 solution}
\renewcommand{\memodate}{September 9, 2026}


%\usepackage[T1]{fontenc}
%\usepackage{lmodern}
%\usepackage{microtype}
%\usepackage[margin=1.25in]{geometry}
%\usepackage{mathtools,amssymb}
\usepackage{xcolor}
%\usepackage[colorlinks=true,linkcolor=blue!55!black,citecolor=blue!55!black,urlcolor=blue!55!black]{hyperref}

%\newcommand{\N}{\mathbb N}
%\newcommand{\Z}{\mathbb Z}
\newcommand{\Fq}{\mathds F_q}
\newcommand{\qbinom}[2]{\genfrac{[}{]}{0pt}{}{#1}{#2}_{q}}
\newcommand{\cS}{\mathfrak S}

\title[Positivity for alternating sums of $q$-factorials]
  {Coefficientwise positivity for\protect\\ alternating sums of $q$-factorials}
\author{Claude (Fable-5.1), with ChatGPT-6 (Astra)}
 \thanks{Directed by D.~E.~Anderson, \texttt{anderson.2804@math.osu.edu}}
\date{September 9, 2026}

%\theoremstyle{plain}
%\newtheorem{theorem}{Theorem}[section]
%\newtheorem{proposition}[theorem]{Proposition}
%\newtheorem{lemma}[theorem]{Lemma}
%\newtheorem{corollary}[theorem]{Corollary}
%\theoremstyle{remark}
%\newtheorem{remark}[theorem]{Remark}

\begin{document}

\begin{abstract}
We prove that the alternating sum
$\sum_{i=0}^{k}(-1)^i\binom ki[m-i]!_q$ has nonnegative coefficients
whenever $m\ge 2k-1$. For $k\ge3$, this range is exact.
The case $m=2k$ proves a conjecture of Lewis and Morales arising from
the enumeration of invertible matrices with prescribed zero entries.
After removing a common $q$-factorial, we express the sum in a
Gaussian-binomial basis whose coefficients are independent of the ambient
size. Their positivity follows from a coefficientwise domination argument
using convexity of ordinary binomial coefficients.
\end{abstract}

\maketitle

\section{The conjecture and the normalized sum}

For a nonnegative integer $r$, write
\[
 [r]=[r]_q=1+q+\cdots+q^{r-1},\qquad
 [r]!=[1][2]\cdots[r],
\]
with $[0]=0$ and $[0]!=1$. Throughout, $\N$ includes $0$, so
membership in $\N[q]$ means coefficientwise nonnegativity.

Lewis and Morales conjectured that
\[
 \sum_{i=0}^{n}(-1)^i\binom ni[2n-i]!\in\N[q]
 \qquad(n\ge1).
\]
This is Conjecture~6.8 of \cite{LM}, also known as OPAC-012.
The sum occurs, up to a power of $q$, in the reduced count of invertible
$2n\times2n$ matrices over a finite field with $n$ prescribed zero entries
in distinct rows and columns. This counting interpretation explains its
nonnegative values at prime powers, but does not imply that its
coefficients are nonnegative. We prove the following stronger statement.

\begin{theorem}\label{thm:main}
For integers $0\le k\le m$, set
\begin{equation}\label{eq:S}
 S_{k,m}(q)=\sum_{i=0}^{k}(-1)^i\binom ki[m-i]!.
\end{equation}
If $m\ge2k-1$, then $S_{k,m}(q)\in\N[q]$. For $k\ge3$, this bound is
sharp: $S_{k,m}(q)$ has a negative coefficient whenever
$k\le m\le2k-2$.
In particular, $S_{n,2n}(q)\in\N[q]$ for every $n\ge1$.
\end{theorem}

The normalization that makes the proof work is to remove the common
factor $[m-k]!$. For $a,r\ge0$, put
\[
 \Pi_r(a)=\prod_{s=1}^{r}[a+s],\qquad \Pi_0(a)=1,
\]
and define
\begin{equation}\label{eq:F}
 F_t(a)=\sum_{r=0}^{t}(-1)^{t-r}\binom tr\Pi_r(a)
 \qquad(t,a\ge0).
\end{equation}
Changing the index in \eqref{eq:S} gives
\begin{equation}\label{eq:SF}
 S_{k,m}(q)=[m-k]!\,F_k(m-k).
\end{equation}
Thus the range $m\ge2k-1$ becomes $a\ge k-1$.

The proof now has two steps. In Section~\ref{sec:expansion}, we expand
$F_t(a)$ as a sum of one explicit term, nonnegative for $a\ge t-1$,
and Gaussian binomial coefficients with universal coefficients
$\Theta_{t,j}(q)$. These coefficients do not depend on $a$.
In Section~\ref{sec:positivity}, we prove their nonnegativity by showing
that the leading positive term dominates all the other terms in their
alternating definition. The underlying comparison reduces to convexity
of ordinary binomial coefficients. The proof of positivity is then
complete; Section~\ref{sec:sharpness} establishes sharpness by computing
the first two possible nonzero coefficients, and
Section~\ref{sec:matrices} records the matrix-counting interpretation.

\section{A universal Gaussian-binomial expansion}\label{sec:expansion}

For $a\ge1$ and $j\ge0$, we use the product formula
\[
 \qbinom{a+j-1}{j}
 =\frac{\prod_{s=0}^{j-1}[a+s]}{[j]!}.
\]
For $a=0$, the right side defines the same notation: it is $1$ when
$j=0$ and $0$ otherwise. All these polynomials belong to $\N[q]$.
We also use $[b]=(1-q^b)/(1-q)$ for negative integers $b$; identities
involving these expressions are understood in $\Z[q,q^{-1}]$.

The universal coefficients in our expansion have the following explicit
form. For $t\ge1$, put $N=t-1$ and define, for $0\le j\le N$,
\begin{equation}\label{eq:Theta}
 \Theta_{t,j}(q)
 =[N]![N]+\sum_{d=1}^{N-j}(-1)^{d+1}q^{N-d}[N-d]!
       \sum_{s=d}^{N}\binom{s}{d}q^{N-s}.
\end{equation}
In particular, $\Theta_{1,0}=0$. Their nonnegativity is not apparent
from this alternating sum; it will be proved in the next section.

\begin{proposition}[Universal expansion]\label{prop:decomposition}
For every $t\ge1$ and $a\ge0$,
\begin{equation}\label{eq:positive-decomposition}
\begin{split}
 F_t(a)={}&q^t[a-t+1]\prod_{r=1}^{t-1}[a+r]\\
 &\quad+q^{a+1}\sum_{j=0}^{t-1}q^j
       \qbinom{a+j-1}{j}\Theta_{t,j}(q).
\end{split}
\end{equation}
\end{proposition}

The first summand vanishes at $a=t-1$ and is nonnegative for
$a\ge t-1$. In the remaining sum, all dependence on $a$, apart from
the common power $q^{a+1}$, lies in the Gaussian binomial coefficients.
Consequently, proving $\Theta_{t,j}\in\N[q]$ will establish
$F_t(a)\in\N[q]$ throughout the required range.

\begin{proof}
We first isolate the explicit summand by induction on $t$, and then
expand the remainder in Gaussian binomial coefficients. The two
identities we need are
\begin{align}
 F_{t+1}(a)&=[a+1]F_t(a+1)-F_t(a),\label{eq:Frec}\\
 \Pi_r(a)&=\sum_{j=0}^{r}q^j[r]!\qbinom{a+j-1}{j}.
 \label{eq:basis}
\end{align}
The first follows from Pascal's identity and
$[a+1]\Pi_r(a+1)=\Pi_{r+1}(a)$. For the second, iterate
\[
 \Pi_r(a)=[r]\Pi_{r-1}(a)
       +q^r\prod_{s=0}^{r-1}[a+s],
\]
which follows from $[a+r]=[r]+q^r[a]$.

To keep the induction readable, abbreviate the inner sums in
\eqref{eq:Theta} by
\begin{equation}\label{eq:lambda}
 \lambda_{N,d}=\sum_{s=d}^{N}\binom{s}{d}q^{N-s}
 \qquad(0\le d\le N),
\end{equation}
and set them to zero outside this range. Pascal's identity gives
\begin{equation}\label{eq:lambdarec}
 \lambda_{N+1,d}=\lambda_{N,d}+\lambda_{N,d-1}
 \quad(d\ge1),\qquad \lambda_{N,0}=[N+1].
\end{equation}
Write the proposed explicit summand as
\[
 \mu_t(a)=q^t[a-t+1]\Pi_{t-1}(a).
\]
For the induction, define a remainder $T_t(a)=\sum_r g^{(t)}_r\Pi_r(a)$,
where
\[
 g^{(t)}_{t-1}=[t-1],\qquad
 g^{(t)}_r=(-1)^{t-r}q^r\lambda_{t-1,t-1-r}
 \quad(0\le r\le t-2),
\]
and $g^{(t)}_r=0$ otherwise. We claim that
\begin{equation}\label{eq:FT}
 F_t(a)=\mu_t(a)+q^{a+1}T_t(a).
\end{equation}

The base case is $F_1(a)=q[a]=\mu_1(a)$, with $T_1=0$.
To pass from $t$ to $t+1$, the explicit term satisfies
\begin{equation}\label{eq:murec}
\begin{split}
 [a+1]\mu_t(a+1)-\mu_t(a)
 ={}&\mu_{t+1}(a)\\
 &+q^{a+1}\bigl([2t-1]\Pi_{t-1}(a)+\Pi_t(a)\bigr).
\end{split}
\end{equation}
For completeness, substitution reduces this identity to
\[
 [a-t+2][a+t]-[a-t+1]-q[a-t][a+t]
 =q^{a-t+1}\bigl([2t-1]+[a+t]\bigr),
\]
which follows from $[x]-[y]=q^y[x-y]$.
Meanwhile, the remainder transforms according to
\[
 q[a+1]\Pi_r(a+1)-\Pi_r(a)=q\Pi_{r+1}(a)-\Pi_r(a).
\]
Applying \eqref{eq:Frec} to \eqref{eq:FT}, we therefore obtain the
desired formula at $t+1$ precisely when
\begin{equation}\label{eq:grec}
 g^{(t+1)}_r=[2t-1]\delta_{r,t-1}+\delta_{r,t}
       +qg^{(t)}_{r-1}-g^{(t)}_r,
\end{equation}
where $\delta$ is the Kronecker delta. For $0\le r\le t-2$, this is
\eqref{eq:lambdarec}. At $r=t-1$, it follows from
\[
 [2t-1]-[t-1]=q^{t-1}[t],\qquad
 \lambda_{t,1}=[t]+\lambda_{t-1,1},
\]
and at $r=t$, it is $[t]=1+q[t-1]$. This proves \eqref{eq:FT}.

Finally, insert \eqref{eq:basis} into $T_t(a)$. With $N=t-1$, the
coefficient multiplying $q^j\qbinom{a+j-1}{j}$ is
\[
 \sum_{r=j}^{t-1}[r]!\,g^{(t)}_r
 =[N]![N]+\sum_{d=1}^{N-j}(-1)^{d+1}q^{N-d}[N-d]!\lambda_{N,d}
 =\Theta_{t,j},
\]
where $d=N-r$. Thus \eqref{eq:FT} is exactly
\eqref{eq:positive-decomposition}.
\end{proof}

\section{Positivity of the universal coefficients}\label{sec:positivity}

We now prove that the alternating sum \eqref{eq:Theta} is nonnegative.
The useful comparison is stronger than a pairing of consecutive terms:
its leading term dominates the sum of all the other terms with their
signs removed. To carry out this comparison, fix $N\ge1$ and package
the summands as
\begin{equation}\label{eq:E}
 E_0=[N]![N],\qquad
 E_d=q^{N-d}[N-d]!\lambda_{N,d}\quad(1\le d\le N),
\end{equation}
where $\lambda_{N,d}$ is the positive polynomial in \eqref{eq:lambda}.

\begin{lemma}[Coefficientwise domination]\label{lem:domination}
For $N\ge1$ and $0\le D\le N$,
\[
 E_0-\sum_{d=1}^{D}E_d\in\N[q].
\]
\end{lemma}

\begin{proof}
We remove the common factorial before comparing coefficients. Define
\begin{equation}\label{eq:Lambda}
 \Lambda_0=[N],\qquad
 \Lambda_D=[N-D+1]\Lambda_{D-1}-q^{N-D}\lambda_{N,D}.
\end{equation}
Then induction gives
\begin{equation}\label{eq:ELambda}
 [N-D]!\Lambda_D=E_0-\sum_{d=1}^{D}E_d.
\end{equation}
The point of this normalization is that the low-degree coefficients of
$\Lambda_D$ have a simple form. More precisely, we prove by induction that
\begin{equation}\label{eq:Lambda-positive}
 \Lambda_D=P_D+q^{N-D+1}Q_D,\qquad
 P_D=\sum_{r=0}^{N-D-1}\binom{r+D}{D}q^r,\qquad Q_D\in\N[q].
\end{equation}
Empty sums are zero. The base case is $\Lambda_0=[N]=P_0$, with $Q_0=0$.

Suppose \eqref{eq:Lambda-positive} holds for $D-1$, and put $L=N-D$.
In the recurrence \eqref{eq:Lambda}, the part involving $Q_{D-1}$ is
already nonnegative. The part to be examined is
\[
 [L+1]P_{D-1}-q^L\lambda_{N,D}.
\]
For $0\le r<L$, its coefficient of $q^r$ is
\[
 \sum_{v=0}^{r}\binom{v+D-1}{D-1}=\binom{r+D}{D},
\]
as required by $P_D$. At degree $L+u$, for $0\le u\le L$, the
coefficient of the product is
\[
 \sum_{v=u}^{L}\binom{v+D-1}{D-1}
 =\binom ND-\binom{u+D-1}{D},
\]
whereas the coefficient of $q^L\lambda_{N,D}$ is $\binom{N-u}{D}$.
Their difference is therefore
\begin{equation}\label{eq:cu}
 c_u=\binom ND-\binom{u+D-1}{D}-\binom{N-u}{D}.
\end{equation}
There are no terms above degree $2L$ in this part of the recurrence.

Here $c_0=0$, and convexity gives $c_u\ge0$ for the remaining values
of $u$. Indeed, the sequence $x\mapsto\binom{x}{D}$ is discretely
convex for $x\ge D-1$: for $D\ge2$ its second forward difference is
$\binom{x}{D-2}\ge0$, and for $D=1$ it is linear. Hence
\[
 u\longmapsto\binom{u+D-1}{D}+\binom{N-u}{D}
\]
is convex on the integer interval $0\le u\le L$ and attains a maximum
at an endpoint. Its endpoint values are
\[
 \binom ND\qquad\text{and}\qquad
 \binom{N-1}{D}+1\le\binom ND,
\]
the inequality following from Pascal's identity. This also covers
$L=0$, when both endpoints coincide.

Since $c_0=0$, the terms remaining after $P_D$ are divisible by
$q^{L+1}$. Explicitly,
\[
 Q_D=\sum_{u=1}^{L}c_uq^{u-1}+q[L+1]Q_{D-1}\in\N[q].
\]
This proves \eqref{eq:Lambda-positive}, and \eqref{eq:ELambda}
proves the lemma.
\end{proof}

\begin{corollary}\label{cor:Theta}
For every $t\ge1$ and $0\le j\le t-1$, $\Theta_{t,j}(q)\in\N[q]$.
\end{corollary}

\begin{proof}
The case $t=1$ is $\Theta_{1,0}=0$. For $t\ge2$, put $N=t-1$ and
$D=N-j$. Rewriting \eqref{eq:Theta} gives
\[
 \Theta_{t,j}
 =\left(E_0-\sum_{d=1}^{D}E_d\right)
   +2\sum_{\substack{1\le d\le D\\d\text{ odd}}}E_d.
\]
The first summand is nonnegative by Lemma~\ref{lem:domination}, and
the second is a sum of nonnegative polynomials.
\end{proof}

\begin{proof}[Proof of the positivity assertion in Theorem~\ref{thm:main}]
For $k=0$, $S_{0,m}=[m]!$. For $k\ge1$, put $a=m-k$. The hypothesis
$m\ge2k-1$ gives $a\ge k-1$, so Proposition~\ref{prop:decomposition}
and Corollary~\ref{cor:Theta} imply $F_k(a)\in\N[q]$.
Equation~\eqref{eq:SF} gives $S_{k,m}\in\N[q]$.
\end{proof}

\section{Sharpness of the range}\label{sec:sharpness}

Below the threshold, the obstruction occurs in one of the first two
possible nonzero coefficients. We use $[q^d]G$ to denote the coefficient
of $q^d$ in a polynomial or formal power series $G$.

\begin{lemma}\label{lem:low-coefficients}
Let $t\ge3$ and $0\le a\le t-2$. Then $[q^d]F_t(a)=0$ for
$0\le d\le a$, and
\begin{align*}
 [q^{a+1}]F_t(a)&=(-1)^t,\\
 [q^{a+2}]F_t(a)&=\binom{a+1}{t-1}-(-1)^t(t-1).
\end{align*}
\end{lemma}

\begin{proof}
Work in $\Z[[q]]$. A product of two distinct factors from
$q^{a+1},q^{a+2},\ldots$ has degree at least $2a+3$, so
\[
 \Pi_r(a)=\frac{\prod_{s=1}^{r}(1-q^{a+s})}{(1-q)^r}
 \equiv\frac{1-q^{a+1}[r]}{(1-q)^r}\pmod{q^{2a+3}}.
\]
Substituting in \eqref{eq:F} and applying the ordinary binomial theorem
gives
\begin{equation}\label{eq:sharp-series}
 F_t(a)\equiv
 \frac{q^t}{(1-q)^t}
 -q^{a+1}\frac{q^t-(2q-1)^t}{(1-q)^{t+1}}
 \pmod{q^{2a+3}}.
\end{equation}
Both terms on the right have degree greater than $a$. At the next two
degrees, the contribution of the first term is respectively $0$ and
$\binom{a+1}{t-1}$, while
\[
 \frac{(2q-1)^t}{(1-q)^{t+1}}
 =(-1)^t\bigl(1-(t-1)q+O(q^2)\bigr).
\]
The remaining numerator term $q^t$ contributes only in degree at least
$a+t+1$. This gives the asserted coefficients.
\end{proof}

\begin{proof}[Proof of the sharpness assertion in Theorem~\ref{thm:main}]
Suppose $k\ge3$ and $k\le m\le2k-2$, and put $a=m-k$.
If $k$ is odd, Lemma~\ref{lem:low-coefficients} and \eqref{eq:SF} give
\[
 [q^{a+1}]S_{k,m}=-1.
\]
If $k\ge4$ is even, the coefficient of $q$ in $[a]!$ is
$\max(a-1,0)$. Since $\binom{a+1}{k-1}\in\{0,1\}$, the same lemma
gives
\[
 [q^{a+2}]S_{k,m}
 =\binom{a+1}{k-1}-(k-1)+\max(a-1,0)<0.
\]
Thus there is a negative coefficient in either case.
\end{proof}

\section{Matrices with prescribed zero entries}\label{sec:matrices}

We conclude by explaining the counting problem behind the conjecture.
For a board $B\subseteq\{1,\ldots,m\}^2$, let $\mathfrak m_r(B,q)$
be the number of $m\times m$ matrices over $\Fq$ of rank $r$ whose
support is contained in $B$, and write
\[
 M_r(B,q)=\frac{\mathfrak m_r(B,q)}{(q-1)^r}
\]
for the reduced count. A bar denotes the complement in the square.
Suppose $B$ consists of $k$ cells in distinct rows and columns. Choosing
a rank-$r$ matrix supported on $B$ amounts to choosing $r$ cells and
assigning a nonzero entry to each, so $M_r(B,q)=\binom kr$.
The complement identity of Lewis and Morales
\cite[Corollary~2.2]{LM} then gives
\begin{equation}\label{eq:isolated-board}
 M_m(\overline B,q)=q^{\binom m2-k}S_{k,m}(q).
\end{equation}
For $m\ge3$, the power of $q$ is nonnegative because $k\le m$.
Thus Theorem~\ref{thm:main} proves coefficientwise positivity of this
reduced count when $m\ge2k-1$. The cases $m\le2$ in this range follow
directly; in particular, the count for $m=k=1$ is zero.

For a permutation $w=w_1\cdots w_m\in\cS_m$, the diagram convention
in \cite{LM} is the coinversion diagram
\[
 I_w=\{(i,w_j):i<j,\ w_i<w_j\}.
\]
The permutation
\[
 v_n=(2n-1)(2n)(2n-3)(2n-2)\cdots3\,4\,1\,2
\]
has $n$ diagram cells in distinct rows and columns. Hence
\begin{equation}\label{eq:LM-matrix}
 M_{2n}(\overline{I_{v_n}},q)=q^{2n(n-1)}S_{n,2n}(q),
\end{equation}
which is the family in Conjecture~6.8 of \cite{LM}.

The same argument applies to
$w=w_0s_{i_1}\cdots s_{i_k}\in\cS_m$, where $w_0=m\cdots2\,1$,
$s_i$ swaps adjacent positions, and the distinct indices $i_r$ are
pairwise nonadjacent. Such a permutation is a sequence of decreasingly
ordered blocks, each of size one or two, with every size-two block
increasing. It therefore avoids $123$, and its coinversion diagram
consists of $k$ cells in distinct rows and columns. Since $m\ge2k$,
\eqref{eq:isolated-board} yields
$M_m(\overline{I_w},q)\in\N[q]$. This establishes this family of cases
of the broader $123$-avoiding positivity conjecture
\cite[Conjecture~6.9]{LM}.

At $q=1$, the normalized sum $F_t(a)$ counts injections
$\{1,\ldots,t\}\to\{1,\ldots,t+a\}$ having no fixed points, by
inclusion--exclusion. The expansion \eqref{eq:positive-decomposition}
provides a positive $q$-analogue in the range $a\ge t-1$.
A statistic on these injections, or a direct interpretation of the
universal coefficients $\Theta_{t,j}$, would explain combinatorially
the positivity proved here by coefficient comparison.

\begin{thebibliography}{1}
\bibitem{LM}
J.~B. Lewis and A.~H. Morales,
\emph{Rook theory of the finite general linear group},
Experiment. Math. \textbf{29} (2020), no.~3, 328--346;
\href{https://arxiv.org/abs/1707.08192}{arXiv:1707.08192}.

\bibitem[OPAC] J.~B.~Lewis, {\it Open Problems in Algebraic Combinatorics}, blog maintained by S.~Hopkins, \href{https://realopacblog.wordpress.com/2019/09/30/matrix-counting-over-finite-fields/}{realopacblog.wordpress.com/2019/09/30/matrix-counting-over-finite-fields/} (2019).

\end{thebibliography}

\end{document}
